\section{Gauge freedom}

We have seen that the electromagnetic fields can be expressed in terms of the
scalar and vector potentials as
\[
	E = - \nabla V - \frac{\partial A}{\partial t}, \qquad
	B = \nabla \times A.
\]
These relations determine the fields from the potentials, but they do not
determine the potentials uniquely.

Consider a sufficiently regular scalar function
\[
	\chi (r,t).
\]
From a given pair of potentials \(V\) and \(A\), define a new pair by
\[
	A' = A + \nabla \chi, \qquad V' = V -
	\frac{\partial \chi}{\partial t}.
\]
Let us calculate the magnetic field generated by the transformed vector
potential:
\[
	B' = \nabla \times A' = \nabla \times
	\left( A + \nabla \chi  \right).
\]
Using the identity \( \nabla \times (\nabla \chi) = 0\) we obtain
\[
	B' = \nabla \times A = B.
\]
The electric field is also unchanged:
\[
	E' = -\nabla V' - \frac{\partial A'}{\partial t} =
	-\nabla \left(V - \frac{\partial \chi}{\partial t} \right) -
	\frac{ \partial}{\partial t} \left(A + \nabla \chi \right) =
	- \nabla V + \nabla \frac{\partial \chi}{\partial t} -
	\frac{\partial A}{\partial t} - \frac{\partial}{\partial t} \nabla \chi.
\]
Since spatial and temporal derivatives commute,
\[
	\nabla \frac{\partial \chi}{\partial t} =
	\frac{\partial}{\partial t} \nabla \chi,
\]
the two additional terms cancel, leaving
\[
	E' = -\nabla V - \frac{\partial A}{\partial t} = E.
\]
The transformation
\[
	A' = A + \nabla \chi, \qquad V' = V - \frac{\partial \chi}{\partial t}
\]
is called a \emph{gauge transformation}. Since the function \(\chi(r,t)\)
can be chosen arbitrarily without changing the electric and magnetic fields,
the potentials are said to possess \emph{gauge freedom}.


This generalizes the freedom already encountered for an ordinary scalar
potential. Adding a constant to a potential does not change its gradient;
similarly, adding the gradient of a scalar function to the vector potential
does not change its curl. In the electromagnetic case, however, a
time-dependent change in \(A\) must be accompanied by a corresponding change
in \(V\) so that the electric field also remains unchanged.

Therefore the values of \(V\) and \(A\) at an individual point are not
uniquely determined by the electromagnetic fields. Different pairs of
potentials connected by a gauge transformation represent the same
electric and magnetic fields and hence the same classical electromagnetic
situation. Some quantities constructed from potentials are nevertheless
unchanged by a gauge transformation. For example, consider the circulation of
the vector potential around a closed curve \(C\):
\[
	\oint_C A' \cdot dr = \oint_C A \cdot dr + \oint_C \nabla \chi \cdot dr.
\]
For a single-valued function \(\chi\),
\[
	\oint_C \nabla \chi \cdot d r = 0,
\]
and therefore
\[
	\oint_C A' \cdot dr = \oint_C A \cdot dr.
\]
Using Stokes' theorem, this gauge-invariant circulation is equal to the
magnetic flux enclosed by the curve:
\[
	\oint_C A \cdot dr = \int_S B \cdot dS = \Phi_B.
\]
In classical electromagnetism, the measurable force on a charged particle
depends only on \(E\) and \(B\), so the non-uniqueness of the
potentials may appear to make them just auxiliary mathematical objects.







